\documentclass{beamer}

\usepackage{amsmath,amsfonts,amssymb,bm}
\usepackage{physics}
\usepackage{mathtools}

\usetheme{Madrid}

\title{Fisher Information:\\ From Classical Statistics to Quantum Neural Networks}
\author{Graduate Lecture Notes}
\date{}

\begin{document}

\frame{\titlepage}

%------------------------------------------------
\begin{frame}{Outline}
\tableofcontents
\end{frame}

%------------------------------------------------
\section{Classical Fisher Information}

\begin{frame}{Statistical Model}
Consider a parametric probability distribution:
\[
p(x|\theta), \quad \theta \in \mathbb{R}^d
\]

Goal:
\begin{itemize}
\item Estimate $\theta$ from data
\item Quantify sensitivity of distribution to $\theta$
\end{itemize}

\pause

Key object:
\[
\textbf{Score function:} \quad 
\ell_\theta(x) = \partial_\theta \log p(x|\theta)
\]

\end{frame}

%------------------------------------------------
\begin{frame}{Definition of Fisher Information}
The Fisher Information is defined as:

\[
\mathcal{I}(\theta) 
= \mathbb{E}_{x \sim p(x|\theta)} 
\left[ (\partial_\theta \log p(x|\theta))^2 \right]
\]

\pause

Equivalent form:

\[
\mathcal{I}(\theta)
= - \mathbb{E}\left[\partial_\theta^2 \log p(x|\theta)\right]
\]

\pause

Interpretation:
\begin{itemize}
\item Measures curvature of likelihood
\item Quantifies distinguishability of nearby distributions
\end{itemize}

\end{frame}

%------------------------------------------------
\begin{frame}{Derivation of Equivalent Form}
Start with normalization:
\[
\int p(x|\theta) dx = 1
\]

Differentiate:
\[
\partial_\theta \int p(x|\theta) dx = 0
\]

\[
\int \partial_\theta p(x|\theta) dx = 0
\]

\pause

Using:
\[
\partial_\theta p = p \partial_\theta \log p
\]

we obtain:
\[
\mathbb{E}[\partial_\theta \log p] = 0
\]

\pause

Second derivative leads to:
\[
\mathcal{I}(\theta)
= -\mathbb{E}[\partial_\theta^2 \log p]
\]

\end{frame}

%------------------------------------------------
\begin{frame}{Cramér–Rao Bound}
For an unbiased estimator $\hat{\theta}$:

\[
\mathrm{Var}(\hat{\theta}) \geq \frac{1}{\mathcal{I}(\theta)}
\]

\pause

Implication:
\begin{itemize}
\item Fisher information sets fundamental precision limit
\item High $\mathcal{I}(\theta)$ → better estimability
\end{itemize}

\end{frame}

%------------------------------------------------
\section{Fisher Information in Optimization}

\begin{frame}{Likelihood Optimization}
Maximum likelihood:
\[
\theta^* = \arg\max_\theta \sum_i \log p(x_i|\theta)
\]

\pause

Gradient:
\[
\nabla_\theta L = \sum_i \partial_\theta \log p(x_i|\theta)
\]

\pause

Hessian:
\[
H \approx - \mathcal{I}(\theta)
\]

\end{frame}

%------------------------------------------------
\begin{frame}{Natural Gradient}
Standard gradient descent:
\[
\theta_{t+1} = \theta_t - \eta \nabla_\theta L
\]

\pause

Problem:
\begin{itemize}
\item Not invariant under reparametrization
\end{itemize}

\pause

Natural gradient:
\[
\theta_{t+1} 
= \theta_t - \eta \, \mathcal{I}(\theta)^{-1} \nabla_\theta L
\]

\pause

Interpretation:
\begin{itemize}
\item Geometry-aware update
\item Uses information metric
\end{itemize}

\end{frame}

%------------------------------------------------
\begin{frame}{Information Geometry}
Fisher information defines a Riemannian metric:

\[
ds^2 = \sum_{ij} \mathcal{I}_{ij} d\theta_i d\theta_j
\]

\pause

Interpretation:
\begin{itemize}
\item Parameter space becomes curved manifold
\item Natural gradient = steepest descent on manifold
\end{itemize}

\end{frame}

%------------------------------------------------
\section{Quantum Fisher Information}

\begin{frame}{Quantum States}
Quantum statistical model:
\[
\rho(\theta)
\]

Pure state:
\[
\rho = \ket{\psi(\theta)}\bra{\psi(\theta)}
\]

\pause

Goal:
\begin{itemize}
\item Estimate parameter encoded in quantum state
\end{itemize}

\end{frame}

%------------------------------------------------
\begin{frame}{Quantum Fisher Information Definition}
Defined via Symmetric Logarithmic Derivative (SLD):

\[
\partial_\theta \rho 
= \frac{1}{2}(L_\theta \rho + \rho L_\theta)
\]

\pause

Quantum Fisher Information:

\[
F_Q(\theta) = \mathrm{Tr}[\rho L_\theta^2]
\]

\end{frame}

%------------------------------------------------
\begin{frame}{QFI for Pure States}
For $\ket{\psi(\theta)}$:

\[
F_Q = 4 \left(
\langle \partial_\theta \psi | \partial_\theta \psi \rangle
- |\langle \psi | \partial_\theta \psi \rangle|^2
\right)
\]

\pause

Key property:
\begin{itemize}
\item Measures distinguishability of quantum states
\end{itemize}

\end{frame}

%------------------------------------------------
\begin{frame}{Quantum Cramér–Rao Bound}
\[
\mathrm{Var}(\hat{\theta}) \geq \frac{1}{F_Q}
\]

\pause

Important:
\begin{itemize}
\item $F_Q \geq \mathcal{I}_{\text{classical}}$
\item Quantum measurements can saturate this bound
\end{itemize}

\end{frame}

%------------------------------------------------
\section{QFI in Quantum Neural Networks}

\begin{frame}{Parameterized Quantum Circuits}
Quantum neural network:

\[
\ket{\psi(\theta)} = U(\theta)\ket{0}
\]

\pause

Output expectation:
\[
L(\theta) = \langle \psi(\theta) | O | \psi(\theta) \rangle
\]

\end{frame}

%------------------------------------------------
\begin{frame}{Gradient Structure}
Gradient:

\[
\partial_\theta L 
= 2 \, \mathrm{Re} \left(
\langle \partial_\theta \psi | O | \psi \rangle
\right)
\]

\pause

Relates directly to QFI structure:
\[
F_Q \sim \langle \partial_\theta \psi | \partial_\theta \psi \rangle
\]

\end{frame}

%------------------------------------------------
\begin{frame}{Quantum Natural Gradient}
Update rule:

\[
\theta_{t+1}
= \theta_t - \eta \, F_Q^{-1} \nabla_\theta L
\]

\pause

Advantages:
\begin{itemize}
\item Avoids barren plateaus
\item Geometry-aware optimization
\item Faster convergence
\end{itemize}

\end{frame}

%------------------------------------------------
\begin{frame}{Connection to Fubini–Study Metric}
QFI induces metric:

\[
ds^2 = \frac{1}{4} F_Q d\theta^2
\]

\pause

Equivalent to:
\begin{itemize}
\item Fubini–Study metric on Hilbert space
\end{itemize}

\end{frame}

%------------------------------------------------
\begin{frame}{Example: Single-Qubit Circuit}
\[
\ket{\psi(\theta)} = e^{-i\theta \sigma_y/2}\ket{0}
\]

\pause

Compute:
\[
F_Q = 1
\]

\pause

Implication:
\begin{itemize}
\item Uniform sensitivity
\item No barren plateau
\end{itemize}

\end{frame}

%------------------------------------------------
\begin{frame}{Barren Plateaus and QFI}
Observation:
\begin{itemize}
\item Vanishing gradients $\Rightarrow$ small QFI
\end{itemize}

\pause

Thus:
\[
F_Q \approx 0 \Rightarrow \text{hard training}
\]

\pause

Design principle:
\begin{itemize}
\item Circuits should maintain large QFI
\end{itemize}

\end{frame}

%------------------------------------------------
\begin{frame}{Summary}
\begin{itemize}
\item Fisher information = curvature of likelihood
\item Natural gradient uses Fisher metric
\item Quantum Fisher information generalizes this
\item QFI governs optimization in quantum neural networks
\item Key role in avoiding barren plateaus
\end{itemize}

\end{frame}

%------------------------------------------------
\end{document}


%------------------------------------------------
\section{Parameter-Shift Rule and QFI}

\begin{frame}{Parameterized Quantum Gates}
Consider a unitary with generator $G$:
\[
U(\theta) = e^{-i \theta G}
\]

\pause

State:
\[
\ket{\psi(\theta)} = U(\theta)\ket{0}
\]

\pause

Generator properties:
\begin{itemize}
\item Hermitian: $G = G^\dagger$
\item Often Pauli strings: $G^2 = I$
\end{itemize}

\end{frame}

%------------------------------------------------
\begin{frame}{Parameter-Shift Rule}
For expectation value:
\[
L(\theta) = \langle \psi(\theta) | O | \psi(\theta) \rangle
\]

\pause

Gradient:
\[
\partial_\theta L 
= \frac{1}{2}\left[
L\left(\theta + \frac{\pi}{2}\right) 
- L\left(\theta - \frac{\pi}{2}\right)
\right]
\]

\pause

Key:
\begin{itemize}
\item Exact (no approximation)
\item Uses shifted circuits only
\end{itemize}

\end{frame}

%------------------------------------------------
\begin{frame}{Derivative of Quantum State}
Derivative:
\[
\partial_\theta \ket{\psi(\theta)} = -i G \ket{\psi(\theta)}
\]

\pause

Thus:
\[
\langle \partial_\theta \psi | \partial_\theta \psi \rangle 
= \langle \psi | G^2 | \psi \rangle
\]

\[
\langle \psi | \partial_\theta \psi \rangle 
= -i \langle G \rangle
\]

\end{frame}

%------------------------------------------------
\begin{frame}{QFI from Generator}
Insert into QFI formula:
\[
F_Q = 4 \left( 
\langle G^2 \rangle - \langle G \rangle^2
\right)
\]

\pause

Thus:
\[
F_Q = 4 \, \mathrm{Var}(G)
\]

\pause

Interpretation:
\begin{itemize}
\item QFI = fluctuations of generator
\item Large variance $\Rightarrow$ high sensitivity
\end{itemize}

\end{frame}

%------------------------------------------------
\begin{frame}{Connection to Parameter-Shift}
Parameter-shift evaluates:
\[
\partial_\theta L = \langle \psi | i[G,O] | \psi \rangle
\]

\pause

QFI involves:
\[
\langle G^2 \rangle - \langle G \rangle^2
\]

\pause

Key link:
\begin{itemize}
\item Gradient depends on commutator $[G,O]$
\item QFI depends on variance of $G$
\item Both controlled by generator structure
\end{itemize}

\pause

Conclusion:
\begin{itemize}
\item Parameter-shift samples geometry encoded in QFI
\end{itemize}

\end{frame}

%------------------------------------------------
\begin{frame}{Quantum Geometric Tensor}
Define:
\[
\chi_{ij} = \langle \partial_i \psi | \partial_j \psi \rangle 
- \langle \partial_i \psi | \psi \rangle 
\langle \psi | \partial_j \psi \rangle
\]

\pause

Decomposition:
\[
\chi_{ij} = \frac{1}{4}F_{ij} + i \Omega_{ij}
\]

\pause

\begin{itemize}
\item Real part = QFI
\item Imaginary part = Berry curvature
\end{itemize}

\end{frame}

%------------------------------------------------
\section{Quantum Boltzmann Machines}

\begin{frame}{Classical Boltzmann Machine}
Energy-based model:
\[
p(x) = \frac{1}{Z} e^{-E(x)}
\]

\pause

Training:
\[
\nabla_\theta \log p(x)
= -\langle \partial_\theta E \rangle_{\text{data}}
+ \langle \partial_\theta E \rangle_{\text{model}}
\]

\pause

Contrastive divergence appears naturally

\end{frame}

%------------------------------------------------
\begin{frame}{Quantum Boltzmann Machine}
Quantum version:
\[
\rho(\theta) = \frac{e^{-H(\theta)}}{Z}
\]

\pause

Or variational form:
\[
\ket{\psi(\theta)} = U(\theta)\ket{0}
\]

\pause

Loss:
\[
\mathcal{L} = D_{\mathrm{KL}}(p_{\text{data}} || p_\theta)
\]

\end{frame}

%------------------------------------------------
\begin{frame}{Gradient in VQBM}
Gradient:
\[
\partial_\theta \mathcal{L}
= \langle \partial_\theta H \rangle_{\text{model}} 
- \langle \partial_\theta H \rangle_{\text{data}}
\]

\pause

Using parameter-shift:
\[
\partial_\theta \langle O \rangle 
\rightarrow \text{shifted circuits}
\]

\pause

Thus:
\begin{itemize}
\item Training reduces to expectation estimation
\end{itemize}

\end{frame}

%------------------------------------------------
\begin{frame}{QFI in VQBM Training}
Natural gradient:
\[
\dot{\theta} = - F_Q^{-1} \nabla_\theta \mathcal{L}
\]

\pause

Interpretation:
\begin{itemize}
\item QFI defines learning metric
\item Equivalent to quantum information geometry
\end{itemize}

\pause

Important:
\begin{itemize}
\item Stabilizes training
\item Reduces vanishing gradients
\end{itemize}

\end{frame}

%------------------------------------------------
\begin{frame}{Connection to Contrastive Divergence}
Classical CD:
\[
\nabla_\theta D_{\mathrm{KL}}
\]

\pause

Quantum case:
\[
\nabla_\theta D_{\mathrm{KL}} 
\sim \text{QFI-weighted gradient}
\]

\pause

Thus:
\begin{itemize}
\item CD implicitly uses information geometry
\item QFI provides optimal metric correction
\end{itemize}

\end{frame}

%------------------------------------------------
\section{Quantum Sensing and QFI}

\begin{frame}{Quantum Parameter Estimation}
Parameter encoded in unitary:
\[
U(\theta) = e^{-i \theta H}
\]

\pause

State:
\[
\ket{\psi(\theta)} = U(\theta)\ket{\psi_0}
\]

\end{frame}

%------------------------------------------------
\begin{frame}{QFI for Sensing}
\[
F_Q = 4 (\langle H^2 \rangle - \langle H \rangle^2)
\]

\pause

Thus:
\begin{itemize}
\item QFI depends on energy variance
\end{itemize}

\pause

Maximize:
\[
\mathrm{Var}(H)
\]

\end{frame}

%------------------------------------------------
\begin{frame}{Standard Quantum Limit vs Heisenberg Limit}
For $N$ probes:

\[
F_Q \sim N \quad (\text{SQL})
\]

\[
F_Q \sim N^2 \quad (\text{Heisenberg})
\]

\pause

Achieved via:
\begin{itemize}
\item Entanglement
\end{itemize}

\end{frame}

%------------------------------------------------
\begin{frame}{Role of Entanglement}
Separable states:
\[
F_Q \leq N
\]

\pause

Entangled states:
\[
F_Q \sim N^2
\]

\pause

Thus:
\begin{itemize}
\item QFI quantifies metrological usefulness
\end{itemize}

\end{frame}

%------------------------------------------------
\begin{frame}{Example: Two-Qubit Sensor}
Hamiltonian:
\[
H = \sigma_z^{(1)} + \sigma_z^{(2)}
\]

\pause

Bell state:
\[
\frac{1}{\sqrt{2}}(\ket{00} + \ket{11})
\]

\pause

Gives enhanced QFI:
\[
F_Q = 4N^2
\]

\end{frame}

%------------------------------------------------
\begin{frame}{Optimization of Quantum Sensors}
Goal:
\[
\max_\theta F_Q(\theta)
\]

\pause

Use variational circuit:
\[
\ket{\psi(\theta)}
\]

\pause

Optimization:
\[
\theta^* = \arg\max F_Q(\theta)
\]

\pause

Uses:
\begin{itemize}
\item Parameter-shift gradients
\item Quantum natural gradient
\end{itemize}

\end{frame}

%------------------------------------------------
\begin{frame}{Connection to Quantum Neural Networks}
Same structure:
\[
\ket{\psi(\theta)} = U(\theta)\ket{0}
\]

\pause

\begin{itemize}
\item QNN → optimize loss
\item Sensor → optimize QFI
\end{itemize}

\pause

Unified view:
\begin{itemize}
\item Optimization on quantum manifold
\end{itemize}

\end{frame}

%------------------------------------------------
\begin{frame}{Summary of Advanced Connections}
\begin{itemize}
\item Parameter-shift gives exact gradients
\item QFI = variance of generator
\item Natural gradient uses QFI inverse
\item VQBM training benefits from QFI geometry
\item Quantum sensing = QFI maximization
\end{itemize}

\end{frame}
